\chapter{Protocolo de anotación y verificación humana}

Este anexo documenta el procedimiento con el que se construyeron las etiquetas
del conjunto de datos. Se incluye porque los resultados del capítulo de
implementación dependen de él: cuando el texto afirma que una muestra fue
``verificada por un operador'', el procedimiento que define qué significa eso es
este.

\section{Principios}

El objetivo del protocolo no es que el anotador acierte, sino que dos anotadores
---o la misma persona en días distintos--- elijan la misma clase ante la misma
imagen. La consistencia tiene prioridad sobre el criterio individual: un conjunto
consistente vale más que uno perfecto pero ruidoso, porque el ruido de etiqueta
impone un techo de desempeño que ningún modelo supera.

De ese principio se derivan tres reglas de aplicación obligatoria:

\begin{enumerate}
    \item La etiqueta se asigna por la \textbf{silueta completa} del vehículo, no
        por una parte, por el color ni por la marca. Reconocer un modelo
        comercial no decide la clase.
    \item Una muestra contiene un vehículo principal. Si hay varios, se etiqueta
        el que domina la imagen.
    \item Ante la duda no se adivina: la muestra se \textbf{descarta}. Una imagen
        descartada no cuesta nada; una mal etiquetada enseña al modelo un error.
\end{enumerate}

\section{Participación de los conjuntos}

Ninguna partición se realizó por imagen. Un mismo vehículo aparece en decenas de
fotogramas consecutivos, y distribuirlos al azar entre entrenamiento y prueba
haría que el modelo memorice el vehículo en lugar de generalizar. La partición se
establece por \textbf{sesión de grabación} y por \textbf{bloque horario}, y se
verificó que no exista ninguna sesión compartida entre entrenamiento y prueba: el
número de sesiones compartidas reportado es cero.

Para el conjunto de placas se aplicó el mismo criterio sobre bloques de grabación,
de modo que el conjunto de prueba contenga horas de captura no vistas durante el
entrenamiento.

\section{Etiquetado del tipo de vehículo}

La taxonomía comprende cinco clases. La Tabla~\ref{tab:taxonomia_tipo} resume el
criterio y la pista decisiva de cada una.

\begin{longtable}{
        P{0.15}
        P{0.38}
        P{0.37}
    }
    \caption{Clases de tipo de vehículo y criterio de asignación}
    \label{tab:taxonomia_tipo}\\

    \toprule
    \rowcolor{gray!20}
    \textbf{Clase} &
    \textbf{Qué incluye} &
    \textbf{Pista decisiva} \\
    \midrule
    \endfirsthead

    \multicolumn{3}{c}{\textit{Continuación de la Tabla
    \ref{tab:taxonomia_tipo}}} \\
    \toprule
    \rowcolor{gray!20}
    \textbf{Clase} &
    \textbf{Qué incluye} &
    \textbf{Pista decisiva} \\
    \midrule
    \endhead

    moto & Motocicleta. & Dos ruedas. \\
    \midrule

    automóvil & Sedán de tres volúmenes. &
    Capó, cabina y \textbf{baúl saliente}; techo bajo. \\
    \midrule

    vagoneta & Todoterreno, \textit{hatchback}, rural y van de carga o familiar. &
    Techo alto que llega hasta el portón trasero, \textbf{sin baúl saliente}. \\
    \midrule

    camioneta & Camioneta de caja abierta. &
    \textbf{Caja de carga abierta} atrás, incluso en cabina doble. \\
    \midrule

    minibús & Van o minibús de pasajeros. &
    Van \textbf{con ventanas} laterales traseras y muchas plazas. \\
    \bottomrule
    \multicolumn{3}{r}{\footnotesize\textit{Fuente: Elaboración propia}} \\*
\end{longtable}

\subsection{Árbol de decisión}

El orden de las preguntas es parte del protocolo: si se responden en otro orden,
dos anotadores pueden llegar a clases distintas. Se aplica de arriba hacia abajo y
gana la primera condición que se cumpla.

\begin{enumerate}
    \item ¿Tiene dos ruedas? Es \textbf{moto}.
    \item ¿Tiene caja de carga abierta atrás? Es \textbf{camioneta}.
    \item ¿Es un camión o un bus grande, de chasis pesado y muchas ruedas? Se
        \textbf{descarta}: no pertenece a ninguna clase de la taxonomía.
    \item ¿Es una van? Con ventanas y muchas plazas es \textbf{minibús}; cerrada o
        de uso particular es \textbf{vagoneta}.
    \item ¿Es un auto pequeño? Con baúl saliente y techo bajo es
        \textbf{automóvil}; con el techo continuo hasta el portón, sin baúl, es
        \textbf{vagoneta}.
\end{enumerate}

\subsection{Las dos confusiones reales}

El error del clasificador se concentra en dos pares, y ambos tienen una regla de
corte explícita:

\begin{itemize}
    \item \textbf{automóvil $\leftrightarrow$ vagoneta.} ¿Hay un baúl que
        sobresale? Es automóvil. ¿No lo hay? Es vagoneta. Un \textit{hatchback}
        pequeño no es automóvil: es vagoneta.
    \item \textbf{vagoneta $\leftrightarrow$ minibús.} La vagoneta es un
        todoterreno o una van pequeña de uso particular; el minibús es una van de
        transporte de pasajeros, con fila de ventanas y muchas plazas.
\end{itemize}

\subsection{La clase fusionada}

Hasta el 23 de octubre de 2026 existió una sexta clase, \textit{furgón}, destinada
a la van de carga sin ventanas laterales traseras. Se fusionó con
\textit{vagoneta} porque su única pista distintiva ---la ventana trasera--- no es
visible de noche, con el vehículo de cola o con el portón tapado; el F1 de prueba
de esa clase era $0{.}16$ frente a $0{.}66$ de macro. Los $98$ recortes
etiquetados como furgón se reauditaron individualmente: $23$ se reasignaron a
camioneta, $20$ a vagoneta, $6$ a minibús y $49$ se descartaron por no ser
identificables. La taxonomía publicada tiene cinco clases.

\section{Etiquetado del color}

El color se etiqueta como tono predominante de la carrocería sobre once clases. El
criterio es el tono percibido, no el nivel de brillo: por eso un vehículo gris
claro bajo sol directo y el mismo vehículo a la sombra comparten etiqueta. Esa
ambigüedad es la que explica la confusión medida entre las clases acromáticas y se
declara en el análisis de resultados.

A diferencia del tipo, el aumento de datos del clasificador de color aplica
únicamente variación de exposición y nunca de saturación ni de tono: el tono es la
etiqueta que el modelo debe predecir, de modo que perturbarlo sería perturbar la
respuesta correcta.

Se conservan las once clases, pero el informe distingue las que alcanzan soporte
suficiente en prueba (al menos $40$ ejemplares) de las que no. Promediar clases con
una o dos imágenes produce un macro que no describe el desempeño del modelo.

\section{Localización de placas}

El etiquetado de placas se realizó en tres etapas encadenadas, con intervención
humana en la segunda.

\subsection{Propuesta automática}

Un motor de reconocimiento de texto de terceros actuó como
\textbf{pseudo-etiquetador} en el sentido estricto del término: propuso una caja
delimitadora por recorte de vehículo, sin que ninguna de esas cajas ingresara al
conjunto de entrenamiento sin revisión. De $4{,}006$ recortes, $727$ ($18{.}1\%$)
contenían una región compatible con una placa; el resto corresponde en su mayoría
a vistas laterales donde la placa no es visible, es decir, negativos reales y no
fallos del etiquetador.

\subsection{Revisión humana}

Un operador revisó $1{,}127$ recortes ---todas las propuestas más una muestra
aleatoria de $400$ negativos--- y aceptó, corrigió, redibujó o descartó cada caja.
El resultado fueron $868$ recortes con placa ($77\%$) y $259$ sin placa
($23\%$). De las $700$ propuestas que sí correspondían a una placa, solo $183$
($26{.}1\%$) resultaron geométricamente utilizables sin corrección: el
etiquetador encontraba la placa, pero su caja recortaba los glifos.

Una caja rechazada no se descarta del todo: se conserva como \textbf{negativo},
que es lo que impide que el detector aprenda a alucinar una placa sobre cualquier
vehículo. La participación resultante fue de $768$ recortes de entrenamiento, $181$
de validación y $178$ de prueba.

\section{Verificación del texto de matrícula}

La revisión de cajas responde \emph{dónde} está la placa; el reconocedor responde
\emph{qué} dice. Ninguna de las dos mide el sistema ensamblado sobre un recorte de
vehículo real, porque para eso hace falta el texto de una placa que el sistema no
haya visto y la revisión de cajas nunca lo pidió.

Por eso se agregó una tercera etapa: un operador escribe el texto de las placas ya
verificadas, sobre una muestra acotada. Esa muestra no es un conjunto de
entrenamiento, es la regla con la que se mide el sistema completo. De las $100$
placas propuestas para verificación, $88$ resultaron legibles y conformaron el
conjunto de medición de extremo a extremo.

\section{Registro y trazabilidad}

Cada muestra se inscribe en un archivo de manifiesto que conserva su ruta, su
etiqueta y, cuando corresponde, una marca que la identifica como sintética. Esa
marca es la que permite separar en cualquier momento el material real del generado
y auditar qué proporción del entrenamiento provino de cada origen.

Los conteos publicados de cada conjunto son verificables contra el manifiesto y
contra los archivos de imagen, y se resumen en la Tabla~\ref{tab:anexo_conteos}.

\begin{table}[H]
    \centering
    \caption{Composición de los conjuntos construidos}
    \label{tab:anexo_conteos}
    \begin{tabular}{P{0.43}P{0.25}P{0.23}}
        \toprule
        \textbf{Conjunto} & \textbf{Muestras} & \textbf{Origen} \\
        \midrule
        Pool de atributos & 4,006 & 5 clases \\
        Tipo: entrenamiento & 2,783 & Pool congelado \\
        Tipo: validación & 590 & Pool congelado \\
        Tipo: prueba & 633 & Pool congelado \\
        Color: entrenamiento & 2,482 & Pool congelado \\
        Color: validación & 539 & Pool congelado \\
        Color: prueba & 551 & Pool congelado \\
        Recortes de placa verificados & 1,237 & Revisión humana \\
        Placas: entrenamiento & 1,085 & Texto confirmado \\
        Placas: prueba & 152 & Texto confirmado \\
        Placas sintéticas & 360,000 & Formato boliviano \\
        Medición de extremo a extremo & 99 & Texto verificado \\
        \bottomrule
    \end{tabular}
\end{table}

El pool de atributos ($4{,}006$ recortes) y la partición de color ($3{,}572$:
$2{,}482 + 539 + 551$) difieren en $434$ recortes. La diferencia se explica por dos
filtros declarados en el protocolo de construcción del conjunto: $426$ recortes
corresponden a minibuses, excluidos porque su librea hace que el color dominante
dependa del ángulo de captura, y $8$ recortes estaban etiquetados con la clase
residual \textit{otro}, que no se adopta como clase de entrenamiento por falta de
soporte y se descarta con el umbral de $10$ muestras. Ambos conjuntos no son
intercambiables y cualquier comparación entre sus métricas debe hacerse sobre el
mismo subconjunto.

\section{Reproductibilidad del etiquetado}

El protocolo se apoya en dos decisiones que lo hacen repetible. La primera es que
las herramientas de revisión son deterministas sobre la misma entrada: la propuesta
automática, la lista de candidatos y la partición por sesión se regeneran igual
mientras el material de origen no cambie. La segunda es que la partición de
ablación queda fijada por una semilla declarada y un archivo de protocolo que
registra semilla y tamaños, de modo que cualquier medición posterior pueda
repetirse sobre exactamente los mismos subconjuntos.
